%%%%%
%
\section{Moyenne-Echantillonnage d'un param\`etre {\tt double}}

\itemstyle{Op\'erateur~:} {\sf MoyEch\_dbl}
\index{MoyEch\_dbl}

%%%%%OBJET%%%%
\itemstyle{Objet~:}
Moyenne et \'echantillonnage d'un vecteur {\tt double}. La
fr\'equence de sortie doit \^etre diviseur de la 
fr\'equence du param\`etre d'entr\'ee. 

%%%%%ARGUMENTS%%%%
\bigskip
%%%%%ENTREES%%%%
\begin{tabarg}
\multicolumn{3}{l}{\titretab {Entr\'ees}}\\
\hline
$X$ & Vecteur, {\tt double} & Param\`etre d'entr\'ee \\
$R_f$ & Coeff., {\tt double} & Rapport des fr\'equences ($ F_{sortie}\over
F_{entree}$ ) \\
$Decal$ & Coeff., {\tt int} & D\'ecalage (\textit{nombre de points})\\
\end{tabarg}

\medskip
%%%%%SORTIES%%%%
\begin{tabarg}
\multicolumn{3}{l}{\titretab {Sorties}}\\
\hline
$X_m$ & Vecteur, {\tt double} & R\'esultat (m\^eme unit\'e que $X$)\\
\end{tabarg}


%%%%%FORMULE%%%%
\itemstyle{Formule}

%%Moy_dbl
$$
Xm[i] = R_f . \sum_{j = i/R_f + Decal}^{(i+1)/R_f -1 + Decal} X[j]
\qquad i=0,Nbre\_Ecbant(X). R_f
$$


%%%%COMMENTAIRE %%%%
\itemstyle{Commentaire~:}

$R_f$ doit \^etre compris dans l'intervalle $]0, 1]$,
et \^etre l'inverse d'un entier. \\
$Decal$ doit \^etre compris dans l'intervalle $[-1/R_f, 1/R_f]$.\\
Le nombre d'\'echantillons du vecteur de sortie est le produit du
nombre d'\'echantillons du vecteur d'entr\'ee par $R_f$, arrondi
\`a l'entier sup\'erieur.

%%%%%AUTEUR%%%%
\itemstyle{Auteur~:}
CNRM/GMEI/TRAMM

%%%%%REFERENCE%%%%
\itemstyle{R\'ef\'erences~:}
